\section{Evaluasi Akhir dan Integrasi Kompetensi}
Evaluasi akhir bertujuan menilai ketercapaian CPMK secara menyeluruh. Anda dapat menggabungkan semua topik menjadi satu proyek mini, lalu melakukan demo dan refleksi. Tahap ini membantu Anda memahami kekuatan dan area yang masih perlu ditingkatkan.

Refleksi juga penting agar Anda bisa merencanakan belajar lanjutan, misalnya memperdalam framework atau keamanan. Dengan evaluasi yang baik, Anda memiliki gambaran jelas tentang kompetensi yang sudah dikuasai. Ini menjadi bekal untuk proyek nyata di dunia kerja.

\begin{table}[ht]
\centering
\begin{tabular}{|P{4cm}|P{7cm}|}
\hline
\textbf{Komponen Evaluasi} & \textbf{Contoh Bukti} \\
\hline
Proyek Akhir & Aplikasi web terintegrasi \\
\hline
Presentasi & Demo fitur dan alur kerja \\
\hline
Refleksi & Ringkasan tantangan dan solusi \\
\hline
\end{tabular}
\caption{Komponen evaluasi akhir}
\end{table}

